\documentclass[11pt]{article}
\usepackage[a4paper,margin=27mm]{geometry}
\usepackage{amsmath,amssymb,amsthm,xcolor,booktabs,array}
\newcommand{\PROVED}{\textcolor{green!40!black}{\textbf{PROVED}}}
\newcommand{\OPEN}{\textcolor{red!70!black}{\textbf{OPEN}}}
\newcommand{\AUDIT}{\textcolor{orange!70!black}{\textbf{AUDIT}}}
\newtheorem{theorem}{Theorem}
\newtheorem{corollary}{Corollary}
\title{Transfer-coefficient audit of the proposed arithmetic intertwiner (v114)}
\date{13 August 2026}

\begin{document}
\maketitle

\section{The zeroth coefficient}

For $0<a<2T$, on $L^2(I_T)$ let
\[
 J_{a,-}F=\frac{\tau_aE_TF-E_TF}{\sqrt2},
 \qquad
 J_{a,+}F=\frac{\tau_aE_TF+E_TF}{\sqrt2}.
\]
Their exact Gram difference is
\begin{equation}
 J_{a,+}^*J_{a,+}-J_{a,-}^*J_{a,-}=S_a+S_{-a}.                 \tag{1}
\end{equation}
Let $P_T$ be the primitive projection and $Q_T=I-P_T$.  The range of
$Q_T$ is the two-dimensional Tate space.

\begin{theorem}[No arithmetic-only unitary changes the minus port into the
plus port --- \PROVED]
There is no isometry $U$ between arithmetic port spaces such that
\begin{equation}
                         UJ_{a,-}P_T=J_{a,+}P_T                 \tag{2}
\end{equation}
for every primitive source.
\end{theorem}

\begin{proof}
If (2) held, equality of the two pulled-back Grams and (1) would give
\begin{equation}
                         P_T(S_a+S_{-a})P_T=0.                 \tag{3}
\end{equation}
The operator $S_a+S_{-a}$ has infinite rank when $0<a<2T$: choose
infinitely many mutually orthogonal functions supported in sufficiently
small disjoint subintervals of an overlap region; their nonzero translated
images remain linearly independent.

On the other hand, if (3) held, decomposition by $P_T+Q_T=I$ would write
\[
 S_a+S_{-a}=P_T(S_a+S_{-a})Q_T
            +Q_T(S_a+S_{-a})P_T
            +Q_T(S_a+S_{-a})Q_T.
\]
Every term on the right has finite rank because $Q_T$ has rank two.  This
contradicts the infinite rank of the left side.
\end{proof}

\begin{corollary}[Minimal-realization uniqueness cannot prove (8a) ---
\PROVED]
The G39/prime-Julia label realization and the physical plus-delay
realization do not have the same transfer function.  They agree in the
logarithmic defect coefficient, but their zeroth coefficients are
$J_{a,-}P_T$ and $J_{a,+}P_T$, respectively.  Therefore equality of their
Krylov word Grams and minimality cannot produce a unitary intertwiner.
\end{corollary}

\begin{proof}
Uniqueness of minimal conservative realizations applies only after equality
of the complete transfer functions.  Equality of their zeroth coefficients
would be (2), excluded by Theorem 1.
\end{proof}

\section{Consequence for the strict Witt budget}

The results of v109--v112 remain valid:
\begin{enumerate}
\item the all-depth arithmetic word norm tends to
$0.5922965364\ldots<1$;
\item the newborn label-to-delay synthesis is an isometry;
\item the Gamma--Tate old/new cross is uniformly bounded.
\end{enumerate}
What fails is the inference that the Witt word vector is the defect source
of the physical Kato residual.  The arithmetic-only version of equation
(8a) in v113 would entail the forbidden zeroth-order intertwining (2).
It is therefore false as stated.

The only type-correct replacement is a cross-place scattering
\begin{equation}
 \mathscr S_T^{\rm jump}:
 \mathscr K_T^>\oplus\mathscr K_T^<\oplus
 \mathscr K_{\Gamma,T}^{\rm aff}\oplus\mathbb C^2_{\rm Tate}
 \longrightarrow
 \mathscr K_T^>\oplus\mathscr K_T^<\oplus
 \mathscr K_{\Gamma,T}^{\rm aff}\oplus\mathbb C^2_{\rm Tate},            \tag{4}
\end{equation}
whose full source pullback sends the joint minus/Gamma/Tate input to the
authoritative plus/Gamma-head output.  It cannot be a direct sum of
arithmetic channel unitaries, by Theorem 1.

For such a source-defined map the required quantitative assertion is
\begin{equation}
 \left\|\mathscr S_T^{\rm jump}
ight\|_{\rm physical\ pullback}\le1.
                                                                    \tag{5}
\end{equation}
By the exact balanced factorization, (5) is equivalent to
$P_TA_TP_T\ge0$.  Hence merely defining (4) from its desired action, or
declaring it passive because its scalar continuation is unimodular, would
be circular.  A proof must derive (4)--(5) from an independent
cross-place conservation law.

\begin{center}
\begin{tabular}{>{\raggedright\arraybackslash}p{.56\linewidth}
                >{\centering\arraybackslash}p{.15\linewidth}
                >{\raggedright\arraybackslash}p{.19\linewidth}}
\toprule
Statement & Status & Location \\
\midrule
Minus and plus physical ports have different primitive Grams & \PROVED & (1)--(3) \\
Arithmetic-only unitary intertwiner & false & Theorem 1 \\
Minimal-realization proof of v113 (8a) & false & Corollary 2 \\
Strict Witt budget and newborn delay isometry & \PROVED & v109--v112 \\
Independent cross-place conservation proving (4)--(5) & \OPEN & not supplied by rows (a)--(c) \\
Condition $G$, row (d) & \OPEN & equivalent to the physical passivity (5) \\
\bottomrule
\end{tabular}
\end{center}

\end{document}
